\documentclass[11pt]{article}

    \usepackage[breakable]{tcolorbox}
    \usepackage{parskip} % Stop auto-indenting (to mimic markdown behaviour)
    

    % Basic figure setup, for now with no caption control since it's done
    % automatically by Pandoc (which extracts ![](path) syntax from Markdown).
    \usepackage{graphicx}
    % Keep aspect ratio if custom image width or height is specified
    \setkeys{Gin}{keepaspectratio}
    % Maintain compatibility with old templates. Remove in nbconvert 6.0
    \let\Oldincludegraphics\includegraphics
    % Ensure that by default, figures have no caption (until we provide a
    % proper Figure object with a Caption API and a way to capture that
    % in the conversion process - todo).
    \usepackage{caption}
    \DeclareCaptionFormat{nocaption}{}
    \captionsetup{format=nocaption,aboveskip=0pt,belowskip=0pt}

    \usepackage{float}
    \floatplacement{figure}{H} % forces figures to be placed at the correct location
    \usepackage{xcolor} % Allow colors to be defined
    \usepackage{enumerate} % Needed for markdown enumerations to work
    \usepackage{geometry} % Used to adjust the document margins
    \usepackage{amsmath} % Equations
    \usepackage{amssymb} % Equations
    \usepackage{textcomp} % defines textquotesingle
    % Hack from http://tex.stackexchange.com/a/47451/13684:
    \AtBeginDocument{%
        \def\PYZsq{\textquotesingle}% Upright quotes in Pygmentized code
    }
    \usepackage{upquote} % Upright quotes for verbatim code
    \usepackage{eurosym} % defines \euro

    \usepackage{iftex}
    \ifPDFTeX
        \usepackage[T1]{fontenc}
        \IfFileExists{alphabeta.sty}{
              \usepackage{alphabeta}
          }{
              \usepackage[mathletters]{ucs}
              \usepackage[utf8x]{inputenc}
          }
    \else
        \usepackage{fontspec}
        \usepackage{unicode-math}
    \fi

    \usepackage{fancyvrb} % verbatim replacement that allows latex
    \usepackage{grffile} % extends the file name processing of package graphics
                         % to support a larger range
    \makeatletter % fix for old versions of grffile with XeLaTeX
    \@ifpackagelater{grffile}{2019/11/01}
    {
      % Do nothing on new versions
    }
    {
      \def\Gread@@xetex#1{%
        \IfFileExists{"\Gin@base".bb}%
        {\Gread@eps{\Gin@base.bb}}%
        {\Gread@@xetex@aux#1}%
      }
    }
    \makeatother
    \usepackage[Export]{adjustbox} % Used to constrain images to a maximum size
    \adjustboxset{max size={0.9\linewidth}{0.9\paperheight}}

    % The hyperref package gives us a pdf with properly built
    % internal navigation ('pdf bookmarks' for the table of contents,
    % internal cross-reference links, web links for URLs, etc.)
    \usepackage{hyperref}
    % The default LaTeX title has an obnoxious amount of whitespace. By default,
    % titling removes some of it. It also provides customization options.
    \usepackage{titling}
    \usepackage{longtable} % longtable support required by pandoc >1.10
    \usepackage{booktabs}  % table support for pandoc > 1.12.2
    \usepackage{array}     % table support for pandoc >= 2.11.3
    \usepackage{calc}      % table minipage width calculation for pandoc >= 2.11.1
    \usepackage[inline]{enumitem} % IRkernel/repr support (it uses the enumerate* environment)
    \usepackage[normalem]{ulem} % ulem is needed to support strikethroughs (\sout)
                                % normalem makes italics be italics, not underlines
    \usepackage{soul}      % strikethrough (\st) support for pandoc >= 3.0.0
    \usepackage{mathrsfs}
    

    
    % Colors for the hyperref package
    \definecolor{urlcolor}{rgb}{0,.145,.698}
    \definecolor{linkcolor}{rgb}{.71,0.21,0.01}
    \definecolor{citecolor}{rgb}{.12,.54,.11}

    % ANSI colors
    \definecolor{ansi-black}{HTML}{3E424D}
    \definecolor{ansi-black-intense}{HTML}{282C36}
    \definecolor{ansi-red}{HTML}{E75C58}
    \definecolor{ansi-red-intense}{HTML}{B22B31}
    \definecolor{ansi-green}{HTML}{00A250}
    \definecolor{ansi-green-intense}{HTML}{007427}
    \definecolor{ansi-yellow}{HTML}{DDB62B}
    \definecolor{ansi-yellow-intense}{HTML}{B27D12}
    \definecolor{ansi-blue}{HTML}{208FFB}
    \definecolor{ansi-blue-intense}{HTML}{0065CA}
    \definecolor{ansi-magenta}{HTML}{D160C4}
    \definecolor{ansi-magenta-intense}{HTML}{A03196}
    \definecolor{ansi-cyan}{HTML}{60C6C8}
    \definecolor{ansi-cyan-intense}{HTML}{258F8F}
    \definecolor{ansi-white}{HTML}{C5C1B4}
    \definecolor{ansi-white-intense}{HTML}{A1A6B2}
    \definecolor{ansi-default-inverse-fg}{HTML}{FFFFFF}
    \definecolor{ansi-default-inverse-bg}{HTML}{000000}

    % common color for the border for error outputs.
    \definecolor{outerrorbackground}{HTML}{FFDFDF}

    % commands and environments needed by pandoc snippets
    % extracted from the output of `pandoc -s`
    \providecommand{\tightlist}{%
      \setlength{\itemsep}{0pt}\setlength{\parskip}{0pt}}
    \DefineVerbatimEnvironment{Highlighting}{Verbatim}{commandchars=\\\{\}}
    % Add ',fontsize=\small' for more characters per line
    \newenvironment{Shaded}{}{}
    \newcommand{\KeywordTok}[1]{\textcolor[rgb]{0.00,0.44,0.13}{\textbf{{#1}}}}
    \newcommand{\DataTypeTok}[1]{\textcolor[rgb]{0.56,0.13,0.00}{{#1}}}
    \newcommand{\DecValTok}[1]{\textcolor[rgb]{0.25,0.63,0.44}{{#1}}}
    \newcommand{\BaseNTok}[1]{\textcolor[rgb]{0.25,0.63,0.44}{{#1}}}
    \newcommand{\FloatTok}[1]{\textcolor[rgb]{0.25,0.63,0.44}{{#1}}}
    \newcommand{\CharTok}[1]{\textcolor[rgb]{0.25,0.44,0.63}{{#1}}}
    \newcommand{\StringTok}[1]{\textcolor[rgb]{0.25,0.44,0.63}{{#1}}}
    \newcommand{\CommentTok}[1]{\textcolor[rgb]{0.38,0.63,0.69}{\textit{{#1}}}}
    \newcommand{\OtherTok}[1]{\textcolor[rgb]{0.00,0.44,0.13}{{#1}}}
    \newcommand{\AlertTok}[1]{\textcolor[rgb]{1.00,0.00,0.00}{\textbf{{#1}}}}
    \newcommand{\FunctionTok}[1]{\textcolor[rgb]{0.02,0.16,0.49}{{#1}}}
    \newcommand{\RegionMarkerTok}[1]{{#1}}
    \newcommand{\ErrorTok}[1]{\textcolor[rgb]{1.00,0.00,0.00}{\textbf{{#1}}}}
    \newcommand{\NormalTok}[1]{{#1}}

    % Additional commands for more recent versions of Pandoc
    \newcommand{\ConstantTok}[1]{\textcolor[rgb]{0.53,0.00,0.00}{{#1}}}
    \newcommand{\SpecialCharTok}[1]{\textcolor[rgb]{0.25,0.44,0.63}{{#1}}}
    \newcommand{\VerbatimStringTok}[1]{\textcolor[rgb]{0.25,0.44,0.63}{{#1}}}
    \newcommand{\SpecialStringTok}[1]{\textcolor[rgb]{0.73,0.40,0.53}{{#1}}}
    \newcommand{\ImportTok}[1]{{#1}}
    \newcommand{\DocumentationTok}[1]{\textcolor[rgb]{0.73,0.13,0.13}{\textit{{#1}}}}
    \newcommand{\AnnotationTok}[1]{\textcolor[rgb]{0.38,0.63,0.69}{\textbf{\textit{{#1}}}}}
    \newcommand{\CommentVarTok}[1]{\textcolor[rgb]{0.38,0.63,0.69}{\textbf{\textit{{#1}}}}}
    \newcommand{\VariableTok}[1]{\textcolor[rgb]{0.10,0.09,0.49}{{#1}}}
    \newcommand{\ControlFlowTok}[1]{\textcolor[rgb]{0.00,0.44,0.13}{\textbf{{#1}}}}
    \newcommand{\OperatorTok}[1]{\textcolor[rgb]{0.40,0.40,0.40}{{#1}}}
    \newcommand{\BuiltInTok}[1]{{#1}}
    \newcommand{\ExtensionTok}[1]{{#1}}
    \newcommand{\PreprocessorTok}[1]{\textcolor[rgb]{0.74,0.48,0.00}{{#1}}}
    \newcommand{\AttributeTok}[1]{\textcolor[rgb]{0.49,0.56,0.16}{{#1}}}
    \newcommand{\InformationTok}[1]{\textcolor[rgb]{0.38,0.63,0.69}{\textbf{\textit{{#1}}}}}
    \newcommand{\WarningTok}[1]{\textcolor[rgb]{0.38,0.63,0.69}{\textbf{\textit{{#1}}}}}
    \makeatletter
    \newsavebox\pandoc@box
    \newcommand*\pandocbounded[1]{%
      \sbox\pandoc@box{#1}%
      % scaling factors for width and height
      \Gscale@div\@tempa\textheight{\dimexpr\ht\pandoc@box+\dp\pandoc@box\relax}%
      \Gscale@div\@tempb\linewidth{\wd\pandoc@box}%
      % select the smaller of both
      \ifdim\@tempb\p@<\@tempa\p@
        \let\@tempa\@tempb
      \fi
      % scaling accordingly (\@tempa < 1)
      \ifdim\@tempa\p@<\p@
        \scalebox{\@tempa}{\usebox\pandoc@box}%
      % scaling not needed, use as it is
      \else
        \usebox{\pandoc@box}%
      \fi
    }
    \makeatother

    % Define a nice break command that doesn't care if a line doesn't already
    % exist.
    \def\br{\hspace*{\fill} \\* }
    % Math Jax compatibility definitions
    \def\gt{>}
    \def\lt{<}
    \let\Oldtex\TeX
    \let\Oldlatex\LaTeX
    \renewcommand{\TeX}{\textrm{\Oldtex}}
    \renewcommand{\LaTeX}{\textrm{\Oldlatex}}
    % Document parameters
    % Document title
    \title{Análisis del Bosón Z: Comparativa de Distribuciones\\en Generación, Reconstrucción y Datos Experimentales}
    \author{Estudio Computacional de Física de Partículas}
    \date{\today}
    
    
    
    
    
    
    
% Pygments definitions
\makeatletter
\def\PY@reset{\let\PY@it=\relax \let\PY@bf=\relax%
    \let\PY@ul=\relax \let\PY@tc=\relax%
    \let\PY@bc=\relax \let\PY@ff=\relax}
\def\PY@tok#1{\csname PY@tok@#1\endcsname}
\def\PY@toks#1+{\ifx\relax#1\empty\else%
    \PY@tok{#1}\expandafter\PY@toks\fi}
\def\PY@do#1{\PY@bc{\PY@tc{\PY@ul{%
    \PY@it{\PY@bf{\PY@ff{#1}}}}}}}
\def\PY#1#2{\PY@reset\PY@toks#1+\relax+\PY@do{#2}}

\@namedef{PY@tok@w}{\def\PY@tc##1{\textcolor[rgb]{0.73,0.73,0.73}{##1}}}
\@namedef{PY@tok@c}{\let\PY@it=\textit\def\PY@tc##1{\textcolor[rgb]{0.24,0.48,0.48}{##1}}}
\@namedef{PY@tok@cp}{\def\PY@tc##1{\textcolor[rgb]{0.61,0.40,0.00}{##1}}}
\@namedef{PY@tok@k}{\let\PY@bf=\textbf\def\PY@tc##1{\textcolor[rgb]{0.00,0.50,0.00}{##1}}}
\@namedef{PY@tok@kp}{\def\PY@tc##1{\textcolor[rgb]{0.00,0.50,0.00}{##1}}}
\@namedef{PY@tok@kt}{\def\PY@tc##1{\textcolor[rgb]{0.69,0.00,0.25}{##1}}}
\@namedef{PY@tok@o}{\def\PY@tc##1{\textcolor[rgb]{0.40,0.40,0.40}{##1}}}
\@namedef{PY@tok@ow}{\let\PY@bf=\textbf\def\PY@tc##1{\textcolor[rgb]{0.67,0.13,1.00}{##1}}}
\@namedef{PY@tok@nb}{\def\PY@tc##1{\textcolor[rgb]{0.00,0.50,0.00}{##1}}}
\@namedef{PY@tok@nf}{\def\PY@tc##1{\textcolor[rgb]{0.00,0.00,1.00}{##1}}}
\@namedef{PY@tok@nc}{\let\PY@bf=\textbf\def\PY@tc##1{\textcolor[rgb]{0.00,0.00,1.00}{##1}}}
\@namedef{PY@tok@nn}{\let\PY@bf=\textbf\def\PY@tc##1{\textcolor[rgb]{0.00,0.00,1.00}{##1}}}
\@namedef{PY@tok@ne}{\let\PY@bf=\textbf\def\PY@tc##1{\textcolor[rgb]{0.80,0.25,0.22}{##1}}}
\@namedef{PY@tok@nv}{\def\PY@tc##1{\textcolor[rgb]{0.10,0.09,0.49}{##1}}}
\@namedef{PY@tok@no}{\def\PY@tc##1{\textcolor[rgb]{0.53,0.00,0.00}{##1}}}
\@namedef{PY@tok@nl}{\def\PY@tc##1{\textcolor[rgb]{0.46,0.46,0.00}{##1}}}
\@namedef{PY@tok@ni}{\let\PY@bf=\textbf\def\PY@tc##1{\textcolor[rgb]{0.44,0.44,0.44}{##1}}}
\@namedef{PY@tok@na}{\def\PY@tc##1{\textcolor[rgb]{0.41,0.47,0.13}{##1}}}
\@namedef{PY@tok@nt}{\let\PY@bf=\textbf\def\PY@tc##1{\textcolor[rgb]{0.00,0.50,0.00}{##1}}}
\@namedef{PY@tok@nd}{\def\PY@tc##1{\textcolor[rgb]{0.67,0.13,1.00}{##1}}}
\@namedef{PY@tok@s}{\def\PY@tc##1{\textcolor[rgb]{0.73,0.13,0.13}{##1}}}
\@namedef{PY@tok@sd}{\let\PY@it=\textit\def\PY@tc##1{\textcolor[rgb]{0.73,0.13,0.13}{##1}}}
\@namedef{PY@tok@si}{\let\PY@bf=\textbf\def\PY@tc##1{\textcolor[rgb]{0.64,0.35,0.47}{##1}}}
\@namedef{PY@tok@se}{\let\PY@bf=\textbf\def\PY@tc##1{\textcolor[rgb]{0.67,0.36,0.12}{##1}}}
\@namedef{PY@tok@sr}{\def\PY@tc##1{\textcolor[rgb]{0.64,0.35,0.47}{##1}}}
\@namedef{PY@tok@ss}{\def\PY@tc##1{\textcolor[rgb]{0.10,0.09,0.49}{##1}}}
\@namedef{PY@tok@sx}{\def\PY@tc##1{\textcolor[rgb]{0.00,0.50,0.00}{##1}}}
\@namedef{PY@tok@m}{\def\PY@tc##1{\textcolor[rgb]{0.40,0.40,0.40}{##1}}}
\@namedef{PY@tok@gh}{\let\PY@bf=\textbf\def\PY@tc##1{\textcolor[rgb]{0.00,0.00,0.50}{##1}}}
\@namedef{PY@tok@gu}{\let\PY@bf=\textbf\def\PY@tc##1{\textcolor[rgb]{0.50,0.00,0.50}{##1}}}
\@namedef{PY@tok@gd}{\def\PY@tc##1{\textcolor[rgb]{0.63,0.00,0.00}{##1}}}
\@namedef{PY@tok@gi}{\def\PY@tc##1{\textcolor[rgb]{0.00,0.52,0.00}{##1}}}
\@namedef{PY@tok@gr}{\def\PY@tc##1{\textcolor[rgb]{0.89,0.00,0.00}{##1}}}
\@namedef{PY@tok@ge}{\let\PY@it=\textit}
\@namedef{PY@tok@gs}{\let\PY@bf=\textbf}
\@namedef{PY@tok@ges}{\let\PY@bf=\textbf\let\PY@it=\textit}
\@namedef{PY@tok@gp}{\let\PY@bf=\textbf\def\PY@tc##1{\textcolor[rgb]{0.00,0.00,0.50}{##1}}}
\@namedef{PY@tok@go}{\def\PY@tc##1{\textcolor[rgb]{0.44,0.44,0.44}{##1}}}
\@namedef{PY@tok@gt}{\def\PY@tc##1{\textcolor[rgb]{0.00,0.27,0.87}{##1}}}
\@namedef{PY@tok@err}{\def\PY@bc##1{{\setlength{\fboxsep}{\string -\fboxrule}\fcolorbox[rgb]{1.00,0.00,0.00}{1,1,1}{\strut ##1}}}}
\@namedef{PY@tok@kc}{\let\PY@bf=\textbf\def\PY@tc##1{\textcolor[rgb]{0.00,0.50,0.00}{##1}}}
\@namedef{PY@tok@kd}{\let\PY@bf=\textbf\def\PY@tc##1{\textcolor[rgb]{0.00,0.50,0.00}{##1}}}
\@namedef{PY@tok@kn}{\let\PY@bf=\textbf\def\PY@tc##1{\textcolor[rgb]{0.00,0.50,0.00}{##1}}}
\@namedef{PY@tok@kr}{\let\PY@bf=\textbf\def\PY@tc##1{\textcolor[rgb]{0.00,0.50,0.00}{##1}}}
\@namedef{PY@tok@bp}{\def\PY@tc##1{\textcolor[rgb]{0.00,0.50,0.00}{##1}}}
\@namedef{PY@tok@fm}{\def\PY@tc##1{\textcolor[rgb]{0.00,0.00,1.00}{##1}}}
\@namedef{PY@tok@vc}{\def\PY@tc##1{\textcolor[rgb]{0.10,0.09,0.49}{##1}}}
\@namedef{PY@tok@vg}{\def\PY@tc##1{\textcolor[rgb]{0.10,0.09,0.49}{##1}}}
\@namedef{PY@tok@vi}{\def\PY@tc##1{\textcolor[rgb]{0.10,0.09,0.49}{##1}}}
\@namedef{PY@tok@vm}{\def\PY@tc##1{\textcolor[rgb]{0.10,0.09,0.49}{##1}}}
\@namedef{PY@tok@sa}{\def\PY@tc##1{\textcolor[rgb]{0.73,0.13,0.13}{##1}}}
\@namedef{PY@tok@sb}{\def\PY@tc##1{\textcolor[rgb]{0.73,0.13,0.13}{##1}}}
\@namedef{PY@tok@sc}{\def\PY@tc##1{\textcolor[rgb]{0.73,0.13,0.13}{##1}}}
\@namedef{PY@tok@dl}{\def\PY@tc##1{\textcolor[rgb]{0.73,0.13,0.13}{##1}}}
\@namedef{PY@tok@s2}{\def\PY@tc##1{\textcolor[rgb]{0.73,0.13,0.13}{##1}}}
\@namedef{PY@tok@sh}{\def\PY@tc##1{\textcolor[rgb]{0.73,0.13,0.13}{##1}}}
\@namedef{PY@tok@s1}{\def\PY@tc##1{\textcolor[rgb]{0.73,0.13,0.13}{##1}}}
\@namedef{PY@tok@mb}{\def\PY@tc##1{\textcolor[rgb]{0.40,0.40,0.40}{##1}}}
\@namedef{PY@tok@mf}{\def\PY@tc##1{\textcolor[rgb]{0.40,0.40,0.40}{##1}}}
\@namedef{PY@tok@mh}{\def\PY@tc##1{\textcolor[rgb]{0.40,0.40,0.40}{##1}}}
\@namedef{PY@tok@mi}{\def\PY@tc##1{\textcolor[rgb]{0.40,0.40,0.40}{##1}}}
\@namedef{PY@tok@il}{\def\PY@tc##1{\textcolor[rgb]{0.40,0.40,0.40}{##1}}}
\@namedef{PY@tok@mo}{\def\PY@tc##1{\textcolor[rgb]{0.40,0.40,0.40}{##1}}}
\@namedef{PY@tok@ch}{\let\PY@it=\textit\def\PY@tc##1{\textcolor[rgb]{0.24,0.48,0.48}{##1}}}
\@namedef{PY@tok@cm}{\let\PY@it=\textit\def\PY@tc##1{\textcolor[rgb]{0.24,0.48,0.48}{##1}}}
\@namedef{PY@tok@cpf}{\let\PY@it=\textit\def\PY@tc##1{\textcolor[rgb]{0.24,0.48,0.48}{##1}}}
\@namedef{PY@tok@c1}{\let\PY@it=\textit\def\PY@tc##1{\textcolor[rgb]{0.24,0.48,0.48}{##1}}}
\@namedef{PY@tok@cs}{\let\PY@it=\textit\def\PY@tc##1{\textcolor[rgb]{0.24,0.48,0.48}{##1}}}

\def\PYZbs{\char`\\}
\def\PYZus{\char`\_}
\def\PYZob{\char`\{}
\def\PYZcb{\char`\}}
\def\PYZca{\char`\^}
\def\PYZam{\char`\&}
\def\PYZlt{\char`\<}
\def\PYZgt{\char`\>}
\def\PYZsh{\char`\#}
\def\PYZpc{\char`\%}
\def\PYZdl{\char`\$}
\def\PYZhy{\char`\-}
\def\PYZsq{\char`\'}
\def\PYZdq{\char`\"}
\def\PYZti{\char`\~}
% for compatibility with earlier versions
\def\PYZat{@}
\def\PYZlb{[}
\def\PYZrb{]}
\makeatother


    % For linebreaks inside Verbatim environment from package fancyvrb.
    \makeatletter
        \newbox\Wrappedcontinuationbox
        \newbox\Wrappedvisiblespacebox
        \newcommand*\Wrappedvisiblespace {\textcolor{red}{\textvisiblespace}}
        \newcommand*\Wrappedcontinuationsymbol {\textcolor{red}{\llap{\tiny$\m@th\hookrightarrow$}}}
        \newcommand*\Wrappedcontinuationindent {3ex }
        \newcommand*\Wrappedafterbreak {\kern\Wrappedcontinuationindent\copy\Wrappedcontinuationbox}
        % Take advantage of the already applied Pygments mark-up to insert
        % potential linebreaks for TeX processing.
        %        {, <, #, %, $, ' and ": go to next line.
        %        _, }, ^, &, >, - and ~: stay at end of broken line.
        % Use of \textquotesingle for straight quote.
        \newcommand*\Wrappedbreaksatspecials {%
            \def\PYGZus{\discretionary{\char`\_}{\Wrappedafterbreak}{\char`\_}}%
            \def\PYGZob{\discretionary{}{\Wrappedafterbreak\char`\{}{\char`\{}}%
            \def\PYGZcb{\discretionary{\char`\}}{\Wrappedafterbreak}{\char`\}}}%
            \def\PYGZca{\discretionary{\char`\^}{\Wrappedafterbreak}{\char`\^}}%
            \def\PYGZam{\discretionary{\char`\&}{\Wrappedafterbreak}{\char`\&}}%
            \def\PYGZlt{\discretionary{}{\Wrappedafterbreak\char`\<}{\char`\<}}%
            \def\PYGZgt{\discretionary{\char`\>}{\Wrappedafterbreak}{\char`\>}}%
            \def\PYGZsh{\discretionary{}{\Wrappedafterbreak\char`\#}{\char`\#}}%
            \def\PYGZpc{\discretionary{}{\Wrappedafterbreak\char`\%}{\char`\%}}%
            \def\PYGZdl{\discretionary{}{\Wrappedafterbreak\char`\$}{\char`\$}}%
            \def\PYGZhy{\discretionary{\char`\-}{\Wrappedafterbreak}{\char`\-}}%
            \def\PYGZsq{\discretionary{}{\Wrappedafterbreak\textquotesingle}{\textquotesingle}}%
            \def\PYGZdq{\discretionary{}{\Wrappedafterbreak\char`\"}{\char`\"}}%
            \def\PYGZti{\discretionary{\char`\~}{\Wrappedafterbreak}{\char`\~}}%
        }
        % Some characters . , ; ? ! / are not pygmentized.
        % This macro makes them "active" and they will insert potential linebreaks
        \newcommand*\Wrappedbreaksatpunct {%
            \lccode`\~`\.\lowercase{\def~}{\discretionary{\hbox{\char`\.}}{\Wrappedafterbreak}{\hbox{\char`\.}}}%
            \lccode`\~`\,\lowercase{\def~}{\discretionary{\hbox{\char`\,}}{\Wrappedafterbreak}{\hbox{\char`\,}}}%
            \lccode`\~`\;\lowercase{\def~}{\discretionary{\hbox{\char`\;}}{\Wrappedafterbreak}{\hbox{\char`\;}}}%
            \lccode`\~`\:\lowercase{\def~}{\discretionary{\hbox{\char`\:}}{\Wrappedafterbreak}{\hbox{\char`\:}}}%
            \lccode`\~`\?\lowercase{\def~}{\discretionary{\hbox{\char`\?}}{\Wrappedafterbreak}{\hbox{\char`\?}}}%
            \lccode`\~`\!\lowercase{\def~}{\discretionary{\hbox{\char`\!}}{\Wrappedafterbreak}{\hbox{\char`\!}}}%
            \lccode`\~`\/\lowercase{\def~}{\discretionary{\hbox{\char`\/}}{\Wrappedafterbreak}{\hbox{\char`\/}}}%
            \catcode`\.\active
            \catcode`\,\active
            \catcode`\;\active
            \catcode`\:\active
            \catcode`\?\active
            \catcode`\!\active
            \catcode`\/\active
            \lccode`\~`\~
        }
    \makeatother

    \let\OriginalVerbatim=\Verbatim
    \makeatletter
    \renewcommand{\Verbatim}[1][1]{%
        %\parskip\z@skip
        \sbox\Wrappedcontinuationbox {\Wrappedcontinuationsymbol}%
        \sbox\Wrappedvisiblespacebox {\FV@SetupFont\Wrappedvisiblespace}%
        \def\FancyVerbFormatLine ##1{\hsize\linewidth
            \vtop{\raggedright\hyphenpenalty\z@\exhyphenpenalty\z@
                \doublehyphendemerits\z@\finalhyphendemerits\z@
                \strut ##1\strut}%
        }%
        % If the linebreak is at a space, the latter will be displayed as visible
        % space at end of first line, and a continuation symbol starts next line.
        % Stretch/shrink are however usually zero for typewriter font.
        \def\FV@Space {%
            \nobreak\hskip\z@ plus\fontdimen3\font minus\fontdimen4\font
            \discretionary{\copy\Wrappedvisiblespacebox}{\Wrappedafterbreak}
            {\kern\fontdimen2\font}%
        }%

        % Allow breaks at special characters using \PYG... macros.
        \Wrappedbreaksatspecials
        % Breaks at punctuation characters . , ; ? ! and / need catcode=\active
        \OriginalVerbatim[#1,codes*=\Wrappedbreaksatpunct]%
    }
    \makeatother

    % Exact colors from NB
    \definecolor{incolor}{HTML}{303F9F}
    \definecolor{outcolor}{HTML}{D84315}
    \definecolor{cellborder}{HTML}{CFCFCF}
    \definecolor{cellbackground}{HTML}{F7F7F7}

    % prompt
    \makeatletter
    \newcommand{\boxspacing}{\kern\kvtcb@left@rule\kern\kvtcb@boxsep}
    \makeatother
    \newcommand{\prompt}[4]{
        {\ttfamily\llap{{\color{#2}[#3]:\hspace{3pt}#4}}\vspace{-\baselineskip}}
    }
    

    
    % Prevent overflowing lines due to hard-to-break entities
    \sloppy
    % Setup hyperref package
    \hypersetup{
      breaklinks=true,  % so long urls are correctly broken across lines
      colorlinks=true,
      urlcolor=urlcolor,
      linkcolor=black,
      citecolor=citecolor,
      }
    % Slightly bigger margins than the latex defaults
    
    \geometry{verbose,tmargin=1in,bmargin=1in,lmargin=1in,rmargin=1in}
    
    

\begin{document}

% ===================================
% PORTADA PROFESIONAL
% ===================================
\begin{titlepage}
    \centering
    \vspace*{\fill}
    
    {\Huge \textbf{Análisis del Bosón Z}} \\
    \vspace{0.5cm}
    {\Large \textbf{Comparativa de Distribuciones en Generación, Reconstrucción y Datos Experimentales}}
    
    \vspace{2cm}
    
    {\large Estudio Computacional de Física de Partículas}
    
    \vspace{3cm}
    
    {\normalsize 
    Este análisis presenta un estudio exhaustivo del bosón Z mediante la comparación de:
    \begin{itemize}
        \item Distribuciones a nivel de verdad de Monte Carlo (generación)
        \item Distribuciones después de la simulación del detector (reconstrucción)
        \item Datos experimentales reales
    \end{itemize}
    
    Se examinan los canales de desintegración en pares electrónicos ($e^-e^+$) y muónicos ($\mu^-\mu^+$)
    para validar la respuesta instrumental del detector y medir la masa invariante del bosón.
    }
    
    \vspace{2cm}
    
    {\normalsize Fecha: \today}
    
    \vspace*{\fill}
\end{titlepage}

% ===================================
% ÍNDICE DE CONTENIDOS
% ===================================
\newpage
\tableofcontents
\newpage

% ===================================
% CONTENIDO PRINCIPAL
% ===================================
    
    Todas las distribuciones que se comparan en el cuaderno están
normalizadas. Eso significa que las figuras se usan para comparar la
forma de las distribuciones y no la cantidad total de eventos, de modo
que la interpretación se centra en cambios de ancho, posición, simetría
y colas entre muestras, canales y niveles de procesamiento.

El objetivo es comparar la cinemática de los leptones y la distribución
de masa invariante en tres etapas del estudio: \textbf{generación
(gen)}, \textbf{reconstrucción (reco)} y \textbf{datos experimentales
(data)}. Esa comparación permite valorar la resolución del detector, el
efecto de la selección de eventos y la compatibilidad entre simulación y
datos.

\hypertarget{introducciuxf3n}{%
\section{Introducción}\label{introducciuxf3n}}

El bosón \(Z\) es una de las partículas mediadoras de la interacción
electrodébil. Su desintegración en pares leptónicos es especialmente
útil porque genera señales limpias y fáciles de reconstruir. Por eso,
los canales \(e^-e^+\) y \(\mu^-\mu^+\) son una buena referencia para
estudiar el comportamiento del detector y medir una masa invariante
cercana al valor esperado del bosón.

En este tipo de análisis, las regiones más estables suelen ser las
centrales del detector, donde la geometría es más uniforme y la
cobertura de los subsistemas de seguimiento y calorimetría es mejor. En
pseudorapidez, esto se traduce en valores cercanos a \(\eta \approx 0\),
mientras que en regiones más adelantadas o más cercanas a los bordes de
aceptación el rendimiento empeora por pérdidas geométricas, menor
resolución y menor eficiencia de reconstrucción. También influyen los
umbrales de disparo en \(p_T\), que afectan sobre todo a eventos menos
energéticos.

    \begin{tcolorbox}[breakable, size=fbox, boxrule=1pt, pad at break*=1mm,colback=cellbackground, colframe=cellborder]
\prompt{In}{incolor}{1}{\boxspacing}
\begin{Verbatim}[commandchars=\\\{\}]
\PY{k+kn}{import}\PY{+w}{ }\PY{n+nn}{ROOT}\PY{+w}{ }\PY{k}{as}\PY{+w}{ }\PY{n+nn}{rt}
\PY{k+kn}{from}\PY{+w}{ }\PY{n+nn}{construccion}\PY{+w}{ }\PY{k+kn}{import} \PY{n}{construccion\PYZus{}individually}
\PY{k+kn}{from}\PY{+w}{ }\PY{n+nn}{total\PYZus{}comparation}\PY{+w}{ }\PY{k+kn}{import} \PY{n}{total\PYZus{}comparation}

\PY{n}{rt}\PY{o}{.}\PY{n}{gROOT}\PY{o}{.}\PY{n}{SetBatch}\PY{p}{(}\PY{k+kc}{True}\PY{p}{)}
\PY{n}{rt}\PY{o}{.}\PY{n}{gErrorIgnoreLevel} \PY{o}{=} \PY{n}{rt}\PY{o}{.}\PY{n}{kError}
\PY{o}{\PYZpc{}}\PY{k}{jsroot} off
\end{Verbatim}
\end{tcolorbox}

    \hypertarget{resultados-de-generaciuxf3n-gen}{%
\section{\texorpdfstring{Resultados de generación
(\texttt{gen})}{Resultados de generación (gen)}}\label{resultados-de-generaciuxf3n-gen}}

En esta sección se presenta la simulación a nivel de verdad de Monte
Carlo. Como todavía no actúa la reconstrucción del detector, este nivel
sirve como referencia física ideal para interpretar lo que vendrá
después. La masa invariante debe mostrar un pico limpio y centrado,
mientras que las variables cinemáticas ayudan a comprobar que la muestra
tiene el comportamiento esperado.

A partir de aquí, cada canal se analiza por separado para que sea más
fácil identificar diferencias entre electrones y muones, y después se
comparan ambos con las muestras reconstruidas y reales.

    \hypertarget{canal-e-e}{%
\subsection{\texorpdfstring{Canal
\(e^-e^+\)}{Canal e\^{}-e\^{}+}}\label{canal-e-e}}

Este canal permite seguir el comportamiento del electrón y del positrón
en la simulación ideal. Es una buena primera comprobación porque la
forma de las distribuciones debe ser regular y el pico de masa
invariante debe aparecer en la zona del bosón \(Z\). Cualquier asimetría
importante o dispersión excesiva aquí señalaría un problema en la
construcción de los eventos.

    \hypertarget{quuxe9-muestran-las-figuras}{%
\subsubsection{Qué muestran las
figuras}\label{quuxe9-muestran-las-figuras}}

La celda siguiente genera las distribuciones de \(p_T\), \(\eta\), los
overlays entre leptones y los cocientes bin a bin. En este nivel, lo
esperado es una buena simetría entre ambas partículas del canal y un
pico de masa invariante bien definido. Si el resultado es limpio,
significa que la muestra de generación está correctamente construida y
puede usarse como base de comparación para el resto del análisis.

    \hypertarget{anuxe1lisis-de-los-resultados-en-generaciuxf3n-e-e}{%
\subsubsection{\texorpdfstring{Análisis de los resultados en generación
(\(e^-e^+\))}{Análisis de los resultados en generación (e\^{}-e\^{}+)}}\label{anuxe1lisis-de-los-resultados-en-generaciuxf3n-e-e}}

    \begin{tcolorbox}[breakable, size=fbox, boxrule=1pt, pad at break*=1mm,colback=cellbackground, colframe=cellborder]
\prompt{In}{incolor}{2}{\boxspacing}
\begin{Verbatim}[commandchars=\\\{\}]
\PY{n}{Electron\PYZus{}gen} \PY{o}{=} \PY{n}{construccion\PYZus{}individually}\PY{p}{(}\PY{l+s+s2}{\PYZdq{}}\PY{l+s+s2}{Data/MC\PYZus{}ee\PYZus{}100K\PYZus{}gen.dat}\PY{l+s+s2}{\PYZdq{}}\PY{p}{,} \PY{n+nb}{type}\PY{o}{=}\PY{l+s+s2}{\PYZdq{}}\PY{l+s+s2}{gen}\PY{l+s+s2}{\PYZdq{}}\PY{p}{)}
\PY{n}{Electron\PYZus{}gen}\PY{o}{.}\PY{n}{plotter}\PY{p}{(}\PY{p}{)}
\end{Verbatim}
\end{tcolorbox}

    \begin{center}
    \adjustimage{max size={0.9\linewidth}{0.9\paperheight}}{boson_files/boson_6_0.png}
    \end{center}
    { \hspace*{\fill} \\}
    
    \begin{center}
    \adjustimage{max size={0.9\linewidth}{0.9\paperheight}}{boson_files/boson_6_1.png}
    \end{center}
    { \hspace*{\fill} \\}
    
    Las figuras muestran distribuciones simétricas de \(p_T\) y \(\eta\)
para ambas partículas, lo que valida que la muestra de generación está
bien construida. El pico de masa invariante debe aparecer limpio
alrededor de 91 GeV/c², sin ensanchamiento significativo. Los ratios
entre electrón y positrón deben estar cercanos a 1, confirmando la
equivalencia física entre ambas partículas por simetría CPT. Esta etapa
actúa como referencia ideal: cualquier diferencia que aparezca después
vendrá de la respuesta instrumental del detector.

    \hypertarget{interpretaciuxf3n-del-ajuste-de-masa}{%
\subsubsection{Interpretación del ajuste de
masa}\label{interpretaciuxf3n-del-ajuste-de-masa}}

Los ajustes permiten resumir la forma del pico con distintos modelos
físicos y empíricos. En esta etapa, lo importante no es solo cuál curva
ajusta mejor, sino cómo se sitúa la masa central y qué tan estrecha
resulta la distribución. Un buen ajuste con \(\chi^2/\mathrm{NDF}\)
razonable y una masa cercana al valor esperado indica que la simulación
reproduce correctamente la resonancia del bosón \(Z\).

    \begin{tcolorbox}[breakable, size=fbox, boxrule=1pt, pad at break*=1mm,colback=cellbackground, colframe=cellborder]
\prompt{In}{incolor}{3}{\boxspacing}
\begin{Verbatim}[commandchars=\\\{\}]
\PY{n}{Electron\PYZus{}gen}\PY{o}{.}\PY{n}{plot\PYZus{}mass\PYZus{}fits}\PY{p}{(}\PY{p}{)}
\end{Verbatim}
\end{tcolorbox}

    \begin{Verbatim}[commandchars=\\\{\}]
Ajuste [GAUS   ] en electron-positron  | Chi2/NDF: 15054.54 / 48  = 313.636
Masa invariante ajustada [GAUS   ] en electron-positron           | m\_ll =
90.587 ± 0.012 GeV/c\^{}2
Ajuste [BW     ] en electron-positron  | Chi2/NDF: 4465.53 / 48  = 93.032
Masa invariante ajustada [BW     ] en electron-positron           | m\_ll =
90.754 ± 0.009 GeV/c\^{}2
Ajuste [VOIGT  ] en electron-positron  | Chi2/NDF: 4465.53 / 47  = 95.011
Masa invariante ajustada [VOIGT  ] en electron-positron           | m\_ll =
90.754 ± 0.009 GeV/c\^{}2
Ajuste [CRUIJFF] en electron-positron  | Chi2/NDF:  799.78 / 45  = 17.773
Masa invariante ajustada [CRUIJFF] en electron-positron           | m\_ll =
91.100 ± 0.019 GeV/c\^{}2
Gráfica de ajustes generada exitosamente: gen\_mass\_fits\_electron\_positron.png
    \end{Verbatim}

    \begin{center}
    \adjustimage{max size={0.9\linewidth}{0.9\paperheight}}{boson_files/boson_9_1.png}
    \end{center}
    { \hspace*{\fill} \\}
    
    Los ajustes con Gaussiana, Breit-Wigner, Voigt y Cruijff deben dar
valores de masa muy similares y cercanos a 91.2 GeV/c². El
\(\chi^2/\mathrm{NDF}\) debe ser bajo, indicando que los datos de
generación siguen una distribución simple y bien modelada. La pequeña
diferencia entre modelos refleja que a este nivel no hay efectos de
resolución importantes, por lo que la forma es más cercana a una
resonancia teórica pura.

    \hypertarget{canal-mu-mu}{%
\subsection{\texorpdfstring{Canal
\(\mu^-\mu^+\)}{Canal \textbackslash mu\^{}-\textbackslash mu\^{}+}}\label{canal-mu-mu}}

    El canal muónico suele mostrar una respuesta más estable que el
electrónico porque el momento del muón se reconstruye con gran
precisión. Por eso es un excelente punto de comparación: si los
resultados del canal muónico son más nítidos, se confirma que la
diferencia viene de la instrumentación y no de la física de fondo.

    \begin{tcolorbox}[breakable, size=fbox, boxrule=1pt, pad at break*=1mm,colback=cellbackground, colframe=cellborder]
\prompt{In}{incolor}{4}{\boxspacing}
\begin{Verbatim}[commandchars=\\\{\}]
\PY{n}{Muon\PYZus{}gen} \PY{o}{=} \PY{n}{construccion\PYZus{}individually}\PY{p}{(}\PY{l+s+s2}{\PYZdq{}}\PY{l+s+s2}{Data/MC\PYZus{}mumu\PYZus{}100K\PYZus{}gen.dat}\PY{l+s+s2}{\PYZdq{}}\PY{p}{,} \PY{n+nb}{type}\PY{o}{=}\PY{l+s+s2}{\PYZdq{}}\PY{l+s+s2}{gen}\PY{l+s+s2}{\PYZdq{}}\PY{p}{)}
\PY{n}{Muon\PYZus{}gen}\PY{o}{.}\PY{n}{plotter}\PY{p}{(}\PY{p}{)}
\end{Verbatim}
\end{tcolorbox}

    \begin{center}
    \adjustimage{max size={0.9\linewidth}{0.9\paperheight}}{boson_files/boson_13_0.png}
    \end{center}
    { \hspace*{\fill} \\}
    
    \begin{center}
    \adjustimage{max size={0.9\linewidth}{0.9\paperheight}}{boson_files/boson_13_1.png}
    \end{center}
    { \hspace*{\fill} \\}
    
    \hypertarget{anuxe1lisis-de-los-resultados-en-generaciuxf3n-mu-mu}{%
\subsubsection{\texorpdfstring{Análisis de los resultados en generación
(\(\mu^-\mu^+\))}{Análisis de los resultados en generación (\textbackslash mu\^{}-\textbackslash mu\^{}+)}}\label{anuxe1lisis-de-los-resultados-en-generaciuxf3n-mu-mu}}

    El canal muónico debe mostrar un comportamiento muy parecido al
electrónico, pero con distribuciones de \(p_T\) y \(\eta\) ligeramente
diferentes según la masa y la vida media del muón. La masa invariante
debe centrar igual que en el canal \(ee\), alrededor de 91 GeV/c². Los
ratios entre muón y antimuón también deben rondar 1, confirmando
equivalencia por simetría. La comparación \(ee\) vs \(\mu\mu\) aquí no
debe mostrar diferencias instrumentales porque ambos están al mismo
nivel de verdad Monte Carlo.

    \begin{tcolorbox}[breakable, size=fbox, boxrule=1pt, pad at break*=1mm,colback=cellbackground, colframe=cellborder]
\prompt{In}{incolor}{5}{\boxspacing}
\begin{Verbatim}[commandchars=\\\{\}]
\PY{n}{Muon\PYZus{}gen}\PY{o}{.}\PY{n}{plot\PYZus{}mass\PYZus{}fits}\PY{p}{(}\PY{p}{)}
\end{Verbatim}
\end{tcolorbox}

    \begin{Verbatim}[commandchars=\\\{\}]
Ajuste [GAUS   ] en muon-anti-muon     | Chi2/NDF: 14111.03 / 48  = 293.980
Masa invariante ajustada [GAUS   ] en muon-anti-muon          | m\_ll = 90.924 ±
0.009 GeV/c\^{}2
Ajuste [BW     ] en muon-anti-muon     | Chi2/NDF: 2134.70 / 48  = 44.473
Masa invariante ajustada [BW     ] en muon-anti-muon          | m\_ll = 90.976 ±
0.007 GeV/c\^{}2
Ajuste [VOIGT  ] en muon-anti-muon     | Chi2/NDF: 2134.70 / 47  = 45.419
Masa invariante ajustada [VOIGT  ] en muon-anti-muon          | m\_ll = 90.976 ±
0.007 GeV/c\^{}2
Ajuste [CRUIJFF] en muon-anti-muon     | Chi2/NDF:  951.44 / 45  = 21.143
Masa invariante ajustada [CRUIJFF] en muon-anti-muon          | m\_ll = 91.121 ±
0.016 GeV/c\^{}2
Gráfica de ajustes generada exitosamente: gen\_mass\_fits\_muon\_anti-muon.png
    \end{Verbatim}

    \begin{Verbatim}[commandchars=\\\{\}]
input\_line\_191:2:24: error: redefinition of 'cruijff'
                double cruijff(double *x, double *par) \{
                       \^{}
input\_line\_186:2:24: note: previous definition is here
                double cruijff(double *x, double *par) \{
                       \^{}
    \end{Verbatim}

    \begin{center}
    \adjustimage{max size={0.9\linewidth}{0.9\paperheight}}{boson_files/boson_16_2.png}
    \end{center}
    { \hspace*{\fill} \\}
    
    Los ajustes del canal muónico en generación también deben converger a
\textasciitilde91 GeV/c² con bajo \(\chi^2/\mathrm{NDF}\). Si el pico es
más estrecho que en \(ee\), eso refleja diferencias en cinemática pero
no en efectos instrumentales (que no aparecen a este nivel). Cuando se
comparen después con \texttt{reco}, las diferencias de ancho reales
indicarán las características de resolución de cada subsistema del
detector.

    \hypertarget{resultados-de-reconstrucciuxf3n-reco}{%
\section{\texorpdfstring{Resultados de reconstrucción
(\texttt{reco})}{Resultados de reconstrucción (reco)}}\label{resultados-de-reconstrucciuxf3n-reco}}

    Aquí se analiza la misma física, pero después de aplicar la respuesta
del detector. Este nivel es clave porque muestra cuánto se ensancha el
pico de masa, cuánto cambia la forma de las distribuciones cinemáticas y
qué diferencias aparecen respecto a la simulación ideal. La comparación
con \texttt{gen} permite separar claramente el efecto físico del efecto
instrumental.

    \hypertarget{canal-e-e-en-reco}{%
\subsection{\texorpdfstring{Canal \(e^-e^+\) en
\texttt{reco}}{Canal e\^{}-e\^{}+ en reco}}\label{canal-e-e-en-reco}}

    En esta parte ya aparece el efecto del detector sobre el canal
electrónico. El pico suele hacerse más ancho que en \texttt{gen} y las
colas pueden cambiar ligeramente por resolución energética y selección
de eventos. Si la forma general sigue estando centrada cerca del bosón
\(Z\), entonces la reconstrucción mantiene la física principal aunque
introduzca ensanchamiento.

    \begin{tcolorbox}[breakable, size=fbox, boxrule=1pt, pad at break*=1mm,colback=cellbackground, colframe=cellborder]
\prompt{In}{incolor}{6}{\boxspacing}
\begin{Verbatim}[commandchars=\\\{\}]
\PY{n}{Electron\PYZus{}rec} \PY{o}{=} \PY{n}{construccion\PYZus{}individually}\PY{p}{(}\PY{l+s+s2}{\PYZdq{}}\PY{l+s+s2}{Data/MC\PYZus{}ee\PYZus{}100K\PYZus{}reco.dat}\PY{l+s+s2}{\PYZdq{}}\PY{p}{,} \PY{n+nb}{type}\PY{o}{=}\PY{l+s+s2}{\PYZdq{}}\PY{l+s+s2}{reco}\PY{l+s+s2}{\PYZdq{}}\PY{p}{)}
\PY{n}{Electron\PYZus{}rec}\PY{o}{.}\PY{n}{plotter}\PY{p}{(}\PY{p}{)}
\end{Verbatim}
\end{tcolorbox}

    \begin{center}
    \adjustimage{max size={0.9\linewidth}{0.9\paperheight}}{boson_files/boson_22_0.png}
    \end{center}
    { \hspace*{\fill} \\}
    
    \begin{center}
    \adjustimage{max size={0.9\linewidth}{0.9\paperheight}}{boson_files/boson_22_1.png}
    \end{center}
    { \hspace*{\fill} \\}
    
    \hypertarget{anuxe1lisis-de-los-resultados-en-reconstrucciuxf3n-e-e}{%
\subsubsection{\texorpdfstring{Análisis de los resultados en
reconstrucción
(\(e^-e^+\))}{Análisis de los resultados en reconstrucción (e\^{}-e\^{}+)}}\label{anuxe1lisis-de-los-resultados-en-reconstrucciuxf3n-e-e}}

Aquí comienzan a verse los efectos del detector. El pico de masa
invariante debe estar más ancho que en \texttt{gen}, reflejando la
resolución limitada de energía en el calorímetro electromagnético. Las
distribuciones de \(p_T\) pueden mostrar cortes en valores bajos por
selección de eventos o umbrales de trigger. La comparación de ratios con
respecto a \texttt{gen} permite cuantificar la pérdida de eficiencia en
diferentes regiones cinemáticas, especialmente a bajo \(p_T\) y en los
bordes de pseudorapidez.

    \begin{tcolorbox}[breakable, size=fbox, boxrule=1pt, pad at break*=1mm,colback=cellbackground, colframe=cellborder]
\prompt{In}{incolor}{7}{\boxspacing}
\begin{Verbatim}[commandchars=\\\{\}]
\PY{n}{Electron\PYZus{}rec}\PY{o}{.}\PY{n}{plot\PYZus{}mass\PYZus{}fits}\PY{p}{(}\PY{p}{)}
\end{Verbatim}
\end{tcolorbox}

    \begin{Verbatim}[commandchars=\\\{\}]
Ajuste [GAUS   ] en electron-positron  | Chi2/NDF: 5049.58 / 48  = 105.200
Masa invariante ajustada [GAUS   ] en electron-positron           | m\_ll =
90.629 ± 0.017 GeV/c\^{}2
Ajuste [BW     ] en electron-positron  | Chi2/NDF: 2288.28 / 48  = 47.673
Masa invariante ajustada [BW     ] en electron-positron           | m\_ll =
90.808 ± 0.016 GeV/c\^{}2
Ajuste [VOIGT  ] en electron-positron  | Chi2/NDF: 1492.16 / 47  = 31.748
Masa invariante ajustada [VOIGT  ] en electron-positron           | m\_ll =
90.776 ± 0.016 GeV/c\^{}2
Ajuste [CRUIJFF] en electron-positron  | Chi2/NDF:  122.31 / 45  = 2.718
Masa invariante ajustada [CRUIJFF] en electron-positron           | m\_ll =
91.472 ± 0.039 GeV/c\^{}2
Gráfica de ajustes generada exitosamente: reco\_mass\_fits\_electron\_positron.png
    \end{Verbatim}

    \begin{Verbatim}[commandchars=\\\{\}]
input\_line\_192:2:24: error: redefinition of 'cruijff'
                double cruijff(double *x, double *par) \{
                       \^{}
input\_line\_186:2:24: note: previous definition is here
                double cruijff(double *x, double *par) \{
                       \^{}
    \end{Verbatim}

    \begin{center}
    \adjustimage{max size={0.9\linewidth}{0.9\paperheight}}{boson_files/boson_24_2.png}
    \end{center}
    { \hspace*{\fill} \\}
    
    El pico de masa en reco debe desplazarse poco respecto a \texttt{gen}
(idealmente menos de 1 GeV), pero con \(\chi^2/\mathrm{NDF}\) mayor por
el ensanchamiento instrumental. Si la Gaussiana ajusta mejor que la
Breit-Wigner es señal de que la resolución del detector domina sobre la
anchura natural del bosón. Comparar los valores de masa entre modelos
ayuda a separar la contribución física de la instrumental.

    \hypertarget{canal-mu-mu-en-reco}{%
\subsection{\texorpdfstring{Canal \(\mu^-\mu^+\) en
\texttt{reco}}{Canal \textbackslash mu\^{}-\textbackslash mu\^{}+ en reco}}\label{canal-mu-mu-en-reco}}

    El canal muónico suele conservar mejor la forma del pico porque la
medida del momento es más precisa. Por eso, al comparar con el canal
electrónico, suele verse una distribución más estrecha y más estable.
Esta diferencia es útil para comprobar que el detector responde como se
espera para cada tipo de leptón.

    \begin{tcolorbox}[breakable, size=fbox, boxrule=1pt, pad at break*=1mm,colback=cellbackground, colframe=cellborder]
\prompt{In}{incolor}{8}{\boxspacing}
\begin{Verbatim}[commandchars=\\\{\}]
\PY{n}{Muon\PYZus{}rec} \PY{o}{=} \PY{n}{construccion\PYZus{}individually}\PY{p}{(}\PY{l+s+s2}{\PYZdq{}}\PY{l+s+s2}{Data/MC\PYZus{}mumu\PYZus{}100K\PYZus{}reco.dat}\PY{l+s+s2}{\PYZdq{}}\PY{p}{,} \PY{n+nb}{type}\PY{o}{=}\PY{l+s+s2}{\PYZdq{}}\PY{l+s+s2}{reco}\PY{l+s+s2}{\PYZdq{}}\PY{p}{)}
\PY{n}{Muon\PYZus{}rec}\PY{o}{.}\PY{n}{plotter}\PY{p}{(}\PY{p}{)}
\end{Verbatim}
\end{tcolorbox}

    \begin{center}
    \adjustimage{max size={0.9\linewidth}{0.9\paperheight}}{boson_files/boson_28_0.png}
    \end{center}
    { \hspace*{\fill} \\}
    
    \begin{center}
    \adjustimage{max size={0.9\linewidth}{0.9\paperheight}}{boson_files/boson_28_1.png}
    \end{center}
    { \hspace*{\fill} \\}
    
    \hypertarget{anuxe1lisis-de-los-resultados-en-reconstrucciuxf3n-mu-mu}{%
\subsubsection{\texorpdfstring{Análisis de los resultados en
reconstrucción
(\(\mu^-\mu^+\))}{Análisis de los resultados en reconstrucción (\textbackslash mu\^{}-\textbackslash mu\^{}+)}}\label{anuxe1lisis-de-los-resultados-en-reconstrucciuxf3n-mu-mu}}

El canal muónico debe mostrar un pico más estrecho que el electrónico
incluso después de reconstrucción, porque el sistema de trazas ofrece
mejor resolución de momento que el calorímetro. Los cortes a bajo
\(p_T\) pueden ser menos severos que en \(ee\) si el subsistema de
muones tiene umbrales diferentes. Esta diferencia relativa entre canales
es una buena prueba de que el modelo del detector es consistente.

    \begin{tcolorbox}[breakable, size=fbox, boxrule=1pt, pad at break*=1mm,colback=cellbackground, colframe=cellborder]
\prompt{In}{incolor}{9}{\boxspacing}
\begin{Verbatim}[commandchars=\\\{\}]
\PY{n}{Muon\PYZus{}rec}\PY{o}{.}\PY{n}{plot\PYZus{}mass\PYZus{}fits}\PY{p}{(}\PY{p}{)}
\end{Verbatim}
\end{tcolorbox}

    \begin{Verbatim}[commandchars=\\\{\}]
Ajuste [GAUS   ] en muon-anti-muon     | Chi2/NDF: 6984.51 / 48  = 145.511
Masa invariante ajustada [GAUS   ] en muon-anti-muon          | m\_ll = 90.800 ±
0.012 GeV/c\^{}2
Ajuste [BW     ] en muon-anti-muon     | Chi2/NDF: 1004.11 / 48  = 20.919
Masa invariante ajustada [BW     ] en muon-anti-muon          | m\_ll = 90.826 ±
0.012 GeV/c\^{}2
Ajuste [VOIGT  ] en muon-anti-muon     | Chi2/NDF:  646.58 / 47  = 13.757
Masa invariante ajustada [VOIGT  ] en muon-anti-muon          | m\_ll = 90.827 ±
0.012 GeV/c\^{}2
Ajuste [CRUIJFF] en muon-anti-muon     | Chi2/NDF:  106.49 / 45  = 2.366
Masa invariante ajustada [CRUIJFF] en muon-anti-muon          | m\_ll = 91.038 ±
0.030 GeV/c\^{}2
Gráfica de ajustes generada exitosamente: reco\_mass\_fits\_muon\_anti-muon.png
    \end{Verbatim}

    \begin{Verbatim}[commandchars=\\\{\}]
input\_line\_193:2:24: error: redefinition of 'cruijff'
                double cruijff(double *x, double *par) \{
                       \^{}
input\_line\_186:2:24: note: previous definition is here
                double cruijff(double *x, double *par) \{
                       \^{}
    \end{Verbatim}

    \begin{center}
    \adjustimage{max size={0.9\linewidth}{0.9\paperheight}}{boson_files/boson_30_2.png}
    \end{center}
    { \hspace*{\fill} \\}
    
    El canal muónico debe mostrar un pico más estrecho que el electrónico
incluso después de reconstrucción, porque el sistema de trazas ofrece
mejor resolución de momento que el calorímetro. Los cortes a bajo
\(p_T\) pueden ser menos severos que en \(ee\) si el subsistema de
muones tiene umbrales diferentes. Esta diferencia relativa entre canales
es una buena prueba de que el modelo del detector es consistente.

    \hypertarget{resultados-con-datos-experimentales-data}{%
\section{\texorpdfstring{Resultados con datos experimentales
(\texttt{data})}{Resultados con datos experimentales (data)}}\label{resultados-con-datos-experimentales-data}}

    En esta sección se trabaja con eventos reales. Aquí ya no basta con que
la simulación sea razonable: los resultados deben ser compatibles con lo
que el detector observa de verdad. Por eso, esta parte del cuaderno es
la validación más importante del análisis completo.

    \hypertarget{canal-e-e-en-data}{%
\subsection{\texorpdfstring{Canal \(e^-e^+\) en
\texttt{data}}{Canal e\^{}-e\^{}+ en data}}\label{canal-e-e-en-data}}

    En datos reales, el canal electrónico permite comprobar si la
reconstrucción y la calibración reproducen bien la masa del bosón \(Z\).
Si el pico aparece en una posición parecida a la de \texttt{reco} y la
forma general se mantiene, la simulación está describiendo correctamente
la respuesta del detector.

    \begin{tcolorbox}[breakable, size=fbox, boxrule=1pt, pad at break*=1mm,colback=cellbackground, colframe=cellborder]
\prompt{In}{incolor}{10}{\boxspacing}
\begin{Verbatim}[commandchars=\\\{\}]
\PY{n}{Electron\PYZus{}data} \PY{o}{=} \PY{n}{construccion\PYZus{}individually}\PY{p}{(}\PY{l+s+s2}{\PYZdq{}}\PY{l+s+s2}{Data/Data\PYZus{}ee\PYZus{}100K\PYZus{}reco.dat}\PY{l+s+s2}{\PYZdq{}}\PY{p}{,} \PY{n+nb}{type}\PY{o}{=}\PY{l+s+s2}{\PYZdq{}}\PY{l+s+s2}{data}\PY{l+s+s2}{\PYZdq{}}\PY{p}{)}
\PY{n}{Electron\PYZus{}data}\PY{o}{.}\PY{n}{plotter}\PY{p}{(}\PY{p}{)}
\end{Verbatim}
\end{tcolorbox}

    \begin{center}
    \adjustimage{max size={0.9\linewidth}{0.9\paperheight}}{boson_files/boson_36_0.png}
    \end{center}
    { \hspace*{\fill} \\}
    
    \begin{center}
    \adjustimage{max size={0.9\linewidth}{0.9\paperheight}}{boson_files/boson_36_1.png}
    \end{center}
    { \hspace*{\fill} \\}
    
    \hypertarget{anuxe1lisis-de-los-resultados-en-datos-experimentales-e-e}{%
\subsubsection{\texorpdfstring{Análisis de los resultados en datos
experimentales
(\(e^-e^+\))}{Análisis de los resultados en datos experimentales (e\^{}-e\^{}+)}}\label{anuxe1lisis-de-los-resultados-en-datos-experimentales-e-e}}

    En datos reales, el pico debe estar en una posición compatible con
\texttt{reco}, validando que la calibración de energía es correcta. Si
hay un desplazamiento significativo, podría indicar un sesgo de
calibración sistemático. La forma debe ser similar a \texttt{reco},
aunque pueden aparecer pequeñas asimetrías por radiación de fotones que
no aparecía en \texttt{gen}. Los cortes cinemáticos deben reflejar los
mismos umbrales de trigger que en \texttt{reco}.

    \begin{tcolorbox}[breakable, size=fbox, boxrule=1pt, pad at break*=1mm,colback=cellbackground, colframe=cellborder]
\prompt{In}{incolor}{11}{\boxspacing}
\begin{Verbatim}[commandchars=\\\{\}]
\PY{n}{Electron\PYZus{}data}\PY{o}{.}\PY{n}{plot\PYZus{}mass\PYZus{}fits}\PY{p}{(}\PY{p}{)}
\end{Verbatim}
\end{tcolorbox}

    \begin{Verbatim}[commandchars=\\\{\}]
Ajuste [GAUS   ] en electron-positron  | Chi2/NDF:  799.09 / 48  = 16.648
Masa invariante ajustada [GAUS   ] en electron-positron           | m\_ll =
89.289 ± 0.067 GeV/c\^{}2
Ajuste [BW     ] en electron-positron  | Chi2/NDF:  378.69 / 48  = 7.889
Masa invariante ajustada [BW     ] en electron-positron           | m\_ll =
89.713 ± 0.063 GeV/c\^{}2
Ajuste [VOIGT  ] en electron-positron  | Chi2/NDF:  369.41 / 47  = 7.860
Masa invariante ajustada [VOIGT  ] en electron-positron           | m\_ll =
89.687 ± 0.062 GeV/c\^{}2
Ajuste [CRUIJFF] en electron-positron  | Chi2/NDF:   46.85 / 45  = 1.041
Masa invariante ajustada [CRUIJFF] en electron-positron           | m\_ll =
91.193 ± 0.164 GeV/c\^{}2
Gráfica de ajustes generada exitosamente: data\_mass\_fits\_electron\_positron.png
    \end{Verbatim}

    \begin{Verbatim}[commandchars=\\\{\}]
input\_line\_194:2:24: error: redefinition of 'cruijff'
                double cruijff(double *x, double *par) \{
                       \^{}
input\_line\_186:2:24: note: previous definition is here
                double cruijff(double *x, double *par) \{
                       \^{}
    \end{Verbatim}

    \begin{center}
    \adjustimage{max size={0.9\linewidth}{0.9\paperheight}}{boson_files/boson_39_2.png}
    \end{center}
    { \hspace*{\fill} \\}
    
    La masa medida en datos debe rondar 91.2 GeV/c² sin desplazamientos
significativos. Si los ajustes convergen a valores ligeramente
diferentes entre modelos, es una señal de que la forma real del pico
contiene componentes asimétricas o colas que no son perfectamente
capturadas por modelos simples. Esto es normal en datos reales y suele
interpretarse como indicador de buena calidad cuando el pico está bien
centrado.

    \hypertarget{canal-mu-mu-en-data}{%
\subsection{\texorpdfstring{Canal \(\mu^-\mu^+\) en
\texttt{data}}{Canal \textbackslash mu\^{}-\textbackslash mu\^{}+ en data}}\label{canal-mu-mu-en-data}}

    Este bloque comprueba la robustez del análisis en muones sobre eventos
reales. Un comportamiento similar al observado en \texttt{reco} sugiere
que la simulación está bien calibrada. Si además el pico sigue siendo
más estrecho que en el canal electrónico, se confirma la ventaja del
sistema de trazas para este tipo de leptones.

    \begin{tcolorbox}[breakable, size=fbox, boxrule=1pt, pad at break*=1mm,colback=cellbackground, colframe=cellborder]
\prompt{In}{incolor}{12}{\boxspacing}
\begin{Verbatim}[commandchars=\\\{\}]
\PY{n}{Muon\PYZus{}data} \PY{o}{=} \PY{n}{construccion\PYZus{}individually}\PY{p}{(}\PY{l+s+s2}{\PYZdq{}}\PY{l+s+s2}{Data/Data\PYZus{}mumu\PYZus{}100K\PYZus{}reco.dat}\PY{l+s+s2}{\PYZdq{}}\PY{p}{,} \PY{n+nb}{type}\PY{o}{=}\PY{l+s+s2}{\PYZdq{}}\PY{l+s+s2}{data}\PY{l+s+s2}{\PYZdq{}}\PY{p}{)}
\PY{n}{Muon\PYZus{}data}\PY{o}{.}\PY{n}{plotter}\PY{p}{(}\PY{p}{)}
\end{Verbatim}
\end{tcolorbox}

    \begin{center}
    \adjustimage{max size={0.9\linewidth}{0.9\paperheight}}{boson_files/boson_43_0.png}
    \end{center}
    { \hspace*{\fill} \\}
    
    \begin{center}
    \adjustimage{max size={0.9\linewidth}{0.9\paperheight}}{boson_files/boson_43_1.png}
    \end{center}
    { \hspace*{\fill} \\}
    
    \hypertarget{anuxe1lisis-de-los-resultados-en-datos-experimentales-mu-mu}{%
\subsubsection{\texorpdfstring{Análisis de los resultados en datos
experimentales
(\(\mu^-\mu^+\))}{Análisis de los resultados en datos experimentales (\textbackslash mu\^{}-\textbackslash mu\^{}+)}}\label{anuxe1lisis-de-los-resultados-en-datos-experimentales-mu-mu}}

    En el canal muónico con datos reales, el pico debe mantenerse más
estrecho que en \(ee\) y en buena compatibilidad con \texttt{reco},
confirmando que la simulación modeló bien la resolución del
espectrómetro de muones. El comportamiento debe ser estable en todas las
regiones cinemáticas accesibles, reflejando que el detector tiene buena
cobertura y homogeneidad.

    \begin{tcolorbox}[breakable, size=fbox, boxrule=1pt, pad at break*=1mm,colback=cellbackground, colframe=cellborder]
\prompt{In}{incolor}{13}{\boxspacing}
\begin{Verbatim}[commandchars=\\\{\}]
\PY{n}{Muon\PYZus{}data}\PY{o}{.}\PY{n}{plot\PYZus{}mass\PYZus{}fits}\PY{p}{(}\PY{p}{)}
\end{Verbatim}
\end{tcolorbox}

    \begin{Verbatim}[commandchars=\\\{\}]
Ajuste [GAUS   ] en muon-anti-muon     | Chi2/NDF: 2903.71 / 48  = 60.494
Masa invariante ajustada [GAUS   ] en muon-anti-muon          | m\_ll = 90.556 ±
0.024 GeV/c\^{}2
Ajuste [BW     ] en muon-anti-muon     | Chi2/NDF:  559.83 / 48  = 11.663
Masa invariante ajustada [BW     ] en muon-anti-muon          | m\_ll = 90.620 ±
0.022 GeV/c\^{}2
Ajuste [VOIGT  ] en muon-anti-muon     | Chi2/NDF:  549.05 / 47  = 11.682
Masa invariante ajustada [VOIGT  ] en muon-anti-muon          | m\_ll = 90.622 ±
0.022 GeV/c\^{}2
Ajuste [CRUIJFF] en muon-anti-muon     | Chi2/NDF:   46.48 / 45  = 1.033
Masa invariante ajustada [CRUIJFF] en muon-anti-muon          | m\_ll = 90.854 ±
0.054 GeV/c\^{}2
Gráfica de ajustes generada exitosamente: data\_mass\_fits\_muon\_anti-muon.png
    \end{Verbatim}

    \begin{Verbatim}[commandchars=\\\{\}]
input\_line\_195:2:24: error: redefinition of 'cruijff'
                double cruijff(double *x, double *par) \{
                       \^{}
input\_line\_186:2:24: note: previous definition is here
                double cruijff(double *x, double *par) \{
                       \^{}
    \end{Verbatim}

    \begin{center}
    \adjustimage{max size={0.9\linewidth}{0.9\paperheight}}{boson_files/boson_46_2.png}
    \end{center}
    { \hspace*{\fill} \\}
    
    La masa medida en datos de muones debe también rondar 91.2 GeV/c² sin
desplazamientos. Comparando con \texttt{reco}, si el acuerdo es bueno,
la predicción del detector es correcta. Si hay pequeñas diferencias,
pueden deberse a efectos de alineación, calibración magnética o
recombinación de trazas que varían ligeramente entre simulación y
realidad.

    \hypertarget{comparativas-globales-entre-tipos-y-canales}{%
\section{Comparativas globales entre tipos y
canales}\label{comparativas-globales-entre-tipos-y-canales}}

    En esta parte se comparan las figuras guardadas para resumir el
comportamiento del análisis completo. La idea no es solo ver si las
curvas se parecen, sino entender qué nos dicen sobre la resolución del
detector, la estabilidad del modelado Monte Carlo y la compatibilidad
entre simulación y datos reales.

Si las distribuciones de \texttt{gen}, \texttt{reco} y \texttt{data}
conservan formas similares y los picos se desplazan poco, entonces el
flujo de reconstrucción está funcionando de manera coherente.

    \hypertarget{canales-superpuestos-por-tipo}{%
\subsection{Canales superpuestos por
tipo}\label{canales-superpuestos-por-tipo}}

    La superposición de \(ee\) y \(\mu\mu\) dentro de un mismo tipo de
muestra permite comprobar si ambos canales reconstruyen una misma
resonancia con diferencias solo instrumentales. Cuando las formas son
parecidas, el resultado apunta a una descripción física consistente;
cuando un canal se ensancha más que el otro, la diferencia se entiende
como efecto del detector.

    \hypertarget{anuxe1lisis-de-comparativas-globales-sobreposiciuxf3n-de-canales}{%
\subsubsection{Análisis de comparativas globales: sobreposición de
canales}\label{anuxe1lisis-de-comparativas-globales-sobreposiciuxf3n-de-canales}}

    \begin{tcolorbox}[breakable, size=fbox, boxrule=1pt, pad at break*=1mm,colback=cellbackground, colframe=cellborder]
\prompt{In}{incolor}{14}{\boxspacing}
\begin{Verbatim}[commandchars=\\\{\}]
\PY{n}{comparacion} \PY{o}{=} \PY{n}{total\PYZus{}comparation}\PY{p}{(}\PY{p}{)}
\PY{n}{comparacion}\PY{o}{.}\PY{n}{plot\PYZus{}invariant\PYZus{}mass\PYZus{}overlay\PYZus{}channels}\PY{p}{(}\PY{n}{normalize}\PY{o}{=}\PY{k+kc}{True}\PY{p}{)}
\end{Verbatim}
\end{tcolorbox}

    \begin{Verbatim}[commandchars=\\\{\}]
Grafica generada: invariant\_mass\_overlay\_channels\_by\_type\_norm.png
    \end{Verbatim}

    \begin{center}
    \adjustimage{max size={0.9\linewidth}{0.9\paperheight}}{boson_files/boson_53_1.png}
    \end{center}
    { \hspace*{\fill} \\}
    
    Cuando se superponen los picos de \(ee\) y \(\mu\mu\), ambos deben estar
centrados en \textasciitilde91 GeV/c² porque es la masa verdadera del Z.
Las diferencias en ancho reflejan la resolución instrumental de cada
subsistema. El ratio debe rondar 1 si ambos canales fueron registrados
con similar eficiencia. Diferencias significativas pueden indicar
desbalances en la aceptancia o contaminación de fondo diferente entre
canales.

    \hypertarget{resultado-de-la-superposiciuxf3n-por-tipo}{%
\subsection{Resultado de la superposición por
tipo}\label{resultado-de-la-superposiciuxf3n-por-tipo}}

    La comparación entre canales dentro de cada muestra permite verificar
que el comportamiento global del bosón \(Z\) se mantiene al pasar de
electrones a muones. Si las diferencias se concentran en la anchura o en
pequeños cambios de forma, la causa más probable es la resolución
distinta de cada subsistema del detector y no una inconsistencia física
del proceso.

    \hypertarget{tipos-superpuestos-por-canal}{%
\subsection{Tipos superpuestos por
canal}\label{tipos-superpuestos-por-canal}}

    Aquí se sigue un canal concreto y se observa cómo evolucionan sus
distribuciones desde \texttt{gen} hasta \texttt{data}. Este resumen es
especialmente útil para ver en un solo vistazo cuánto ensanchamiento
introduce la reconstrucción y si los datos reales quedan alineados con
esa predicción.

    \hypertarget{anuxe1lisis-de-comparativas-por-tipo-de-procesamiento}{%
\subsubsection{Análisis de comparativas por tipo de
procesamiento}\label{anuxe1lisis-de-comparativas-por-tipo-de-procesamiento}}

    \begin{tcolorbox}[breakable, size=fbox, boxrule=1pt, pad at break*=1mm,colback=cellbackground, colframe=cellborder]
\prompt{In}{incolor}{15}{\boxspacing}
\begin{Verbatim}[commandchars=\\\{\}]
\PY{n}{comparacion} \PY{o}{=} \PY{n}{total\PYZus{}comparation}\PY{p}{(}\PY{p}{)}
\PY{n}{comparacion}\PY{o}{.}\PY{n}{plot\PYZus{}invariant\PYZus{}mass\PYZus{}overlay\PYZus{}types\PYZus{}by\PYZus{}channel}\PY{p}{(}\PY{n}{normalize}\PY{o}{=}\PY{k+kc}{True}\PY{p}{)}
\PY{n}{comparacion}\PY{o}{.}\PY{n}{canvas\PYZus{}invariant\PYZus{}mass\PYZus{}overlay\PYZus{}types\PYZus{}by\PYZus{}channel}\PY{o}{.}\PY{n}{Draw}\PY{p}{(}\PY{p}{)}
\end{Verbatim}
\end{tcolorbox}

    \begin{Verbatim}[commandchars=\\\{\}]
Grafica generada: invariant\_mass\_overlay\_types\_by\_channel\_norm.png
    \end{Verbatim}

    \begin{center}
    \adjustimage{max size={0.9\linewidth}{0.9\paperheight}}{boson_files/boson_60_1.png}
    \end{center}
    { \hspace*{\fill} \\}
    
    Las curvas de \texttt{gen}, \texttt{reco} y \texttt{data} superpuestas
muestran la cascada de efectos: la verdad física se amplía por la
resolución del detector en \texttt{reco}, y \texttt{data} debe coincidir
estrechamente con \texttt{reco}. Si \texttt{data} se desvía
significativamente de \texttt{reco}, hay un problema de simulación o
calibración. Especialmente en los bordes de distribuciones, las
diferencias pueden indicar pérdida de aceptancia no modelada o efectos
radiativo no incluidos en MC.

    \hypertarget{resultado-de-la-superposiciuxf3n-por-canal}{%
\subsection{Resultado de la superposición por
canal}\label{resultado-de-la-superposiciuxf3n-por-canal}}

    Este resumen reúne en una sola vista la evolución desde la verdad Monte
Carlo hasta los datos reales. La lectura correcta consiste en comprobar
si el ensanchamiento esperado aparece de forma progresiva y si el nivel
\texttt{data} se mantiene cerca de \texttt{reco}. Cuando eso ocurre, la
cadena de análisis es consistente y la simulación reproduce bien el
experimento.

    \hypertarget{ratios-y-lectura-fuxedsica}{%
\subsection{Ratios y lectura física}\label{ratios-y-lectura-fuxedsica}}

    Los cocientes bin a bin completan la comparación porque cuantifican las
diferencias locales. Un valor cercano a 1 indica buen acuerdo;
desviaciones sistemáticas marcan regiones donde la respuesta
instrumental, la selección o la calibración pueden estar introduciendo
sesgos. Por eso esta parte es la que mejor resume dónde el modelo
funciona bien y dónde conviene revisar el análisis.

    \hypertarget{anuxe1lisis-de-figuras-comparativas-individuales}{%
\subsubsection{Análisis de figuras comparativas
individuales}\label{anuxe1lisis-de-figuras-comparativas-individuales}}

    \begin{tcolorbox}[breakable, size=fbox, boxrule=1pt, pad at break*=1mm,colback=cellbackground, colframe=cellborder]
\prompt{In}{incolor}{16}{\boxspacing}
\begin{Verbatim}[commandchars=\\\{\}]
\PY{n}{comparacion}\PY{o}{.}\PY{n}{plot\PYZus{}comparative\PYZus{}figures\PYZus{}individually}\PY{p}{(}\PY{p}{)}
\end{Verbatim}
\end{tcolorbox}

    \begin{center}
    \adjustimage{max size={0.9\linewidth}{0.9\paperheight}}{boson_files/boson_67_0.png}
    \end{center}
    { \hspace*{\fill} \\}
    
    \begin{center}
    \adjustimage{max size={0.9\linewidth}{0.9\paperheight}}{boson_files/boson_67_1.png}
    \end{center}
    { \hspace*{\fill} \\}
    
    \begin{center}
    \adjustimage{max size={0.9\linewidth}{0.9\paperheight}}{boson_files/boson_67_2.png}
    \end{center}
    { \hspace*{\fill} \\}
    
    \begin{center}
    \adjustimage{max size={0.9\linewidth}{0.9\paperheight}}{boson_files/boson_67_3.png}
    \end{center}
    { \hspace*{\fill} \\}
    
    \begin{center}
    \adjustimage{max size={0.9\linewidth}{0.9\paperheight}}{boson_files/boson_67_4.png}
    \end{center}
    { \hspace*{\fill} \\}
    
    \begin{center}
    \adjustimage{max size={0.9\linewidth}{0.9\paperheight}}{boson_files/boson_67_5.png}
    \end{center}
    { \hspace*{\fill} \\}
    
    \begin{center}
    \adjustimage{max size={0.9\linewidth}{0.9\paperheight}}{boson_files/boson_67_6.png}
    \end{center}
    { \hspace*{\fill} \\}
    
    \begin{center}
    \adjustimage{max size={0.9\linewidth}{0.9\paperheight}}{boson_files/boson_67_7.png}
    \end{center}
    { \hspace*{\fill} \\}
    
    Estos gráficos de lado a lado permiten ver con detalle cómo cada
variable cinemática se transforma a lo largo de la cadena
\texttt{gen\ →\ reco\ →\ data}. Los cortes cinemáticos claros a bajo
\(p_T\) en \texttt{reco} y \texttt{data} (que no aparecen en
\texttt{gen}) son firmas del trigger y de los umbrales de selección del
detector. Las colas asimétricas en \(\eta\) cuando se va a los bordes
reflejan cambios en eficiencia debidos a geometría imperfecta o
cobertura limitada en pseudorapidez.

    \hypertarget{conclusiones}{%
\section{Conclusiones}\label{conclusiones}}

    El análisis muestra que la estructura física del bosón \(Z\) se conserva
en los tres niveles estudiados, aunque la reconstrucción introduce un
ensanchamiento y pequeñas deformaciones esperables por la respuesta del
detector. El canal muónico tiende a ofrecer una resolución mejor que el
electrónico, lo que encaja con la diferencia entre los subsistemas de
medida.

Las discrepancias más visibles aparecen en las regiones de menor energía
transversal y en los bordes de aceptación geométrica. A bajas energías,
algunos leptones no superan los umbrales de disparo o reconstrucción,
por lo que se pierden eventos y la eficiencia cae. Además, en esas zonas
es más fácil que aparezcan confusiones con fondo o con eventos mal
reconstruidos, lo que altera la forma de las distribuciones y puede
ensanchar el pico de masa invariante.

En conjunto, \texttt{gen}, \texttt{reco} y \texttt{data} forman una
cadena coherente: la simulación ideal fija la referencia, la
reconstrucción incorpora el efecto instrumental y los datos reales
validan hasta qué punto el modelo reproduce el experimento. Las
diferencias encontradas se explican sobre todo por resolución limitada,
aceptación geométrica, pérdidas de eventos a bajo \(p_T\) y posible
contaminación de fondo, y no por una inconsistencia física del proceso
principal.


    % Add a bibliography block to the postdoc
    
    
    
\end{document}
